% SGA 3, Expose IV, English reconstruction.
% Scope: Section 4.4 heading through the proof of Lemma 4.4.4, stopping
% immediately before editor's note 32 and Lemma 4.4.5.
% Authority: Polo--Gille Exp4-13oct24.pdf, component-local pages 35--37
% (pp. 35 and 37 partial; combined-reader pages 223--225; printed pp. 205--213
% under the re-edition's dual foliation).
% Jacob Reinhold's CC BY 4.0 Markdown lineage was used only as a phrasing
% comparison. The Polo--Gille PDF controls prose, formulas, numbering, and notes.

\setcounter{footnote}{29}

\SGARefTarget{sga3:iv:subsection:4:4}\subsection*{4.4. 层范畴的正合性质}
\addcontentsline{toc}{subsection}{\SGARefLink{sga3:iv:subsection:4:4}{4.4}. 层范畴的正合性质}

\medskip
\noindent\SGARefTarget{sga3:iv:theorem:4:4:1}\textbf{定理 4.4.1.}
\begin{enumerate}
  \item[(i)]\SGARefTarget{sga3:iv:theorem:4:4:1:item:i:01}
    \textit{\(\widetilde{\mathscr C}\) 中任意逆向极限均存在,
    且\textquotedblleft 在 \(\widehat{\mathscr C}\) 中计算\textquotedblright;
    即包含函子
    \[
      i:\widetilde{\mathscr C}\longrightarrow\widehat{\mathscr C}
    \]
    与逆向极限可交换. 更明确地说, 若 \((X_\alpha)\) 是层的逆向系, 则预层}
    \[
      \varprojlim_\alpha i(X_\alpha):
      S\longmapsto\varprojlim_\alpha X_\alpha(S)
    \]
    \textit{是层, 且}
    \[
      i\!\left(\varprojlim_\alpha X_\alpha\right)
      =\varprojlim_\alpha i(X_\alpha).
    \]

  \item[(ii)]\SGARefTarget{sga3:iv:theorem:4:4:1:item:ii:01}
    \textit{\(\widetilde{\mathscr C}\) 中任意正向极限均存在.
    若 \((X_\alpha)\) 是层的正向系, 则}
    \[
      \varinjlim_\alpha X_\alpha
      =a\!\left(\varinjlim_\alpha i(X_\alpha)\right),
    \]
    \textit{其中 \(\varinjlim_\alpha i(X_\alpha)\)
    是各 \(i(X_\alpha)\) 的预层正向极限:}
    \[
      \varinjlim_\alpha i(X_\alpha):
      S\longmapsto\varinjlim_\alpha X_\alpha(S).
    \]

  \item[(iii)]\SGARefTarget{sga3:iv:theorem:4:4:1:item:iii:01} \textit{函子}
    \[
      a:\widehat{\mathscr C}\longrightarrow\widetilde{\mathscr C}
    \]
    \textit{与任意正向极限及有限逆向极限可交换.}
\end{enumerate}

断言 \SGARefLink{sga3:iv:theorem:4:4:1:item:i:01}{(i)} 和
\SGARefLink{sga3:iv:theorem:4:4:1:item:ii:01}{(ii)} 形式地由伴随公式
(\SGARefLink{sga3:iv:theorem:4:3:14}{4.3.14}) 得到,
而\footnote{编者注: 这里修改了原文.}
\SGARefLink{sga3:iv:theorem:4:4:1:item:iii:01}{(iii)} 的第一个断言由
\SGARefLink{sga3:iv:theorem:4:4:1:item:ii:01}{(ii)} 得到.
\SGARefLink{sga3:iv:theorem:4:4:1:item:iii:01}{(iii)} 的第二个断言
已在 \SGARefLink{sga3:iv:remark:4:3:16}{4.3.16} 中指出.

\medskip
\noindent\SGARefTarget{sga3:iv:scholium:4:4:2}\textbf{附论 4.4.2.}
\textit{借助\SGARefLink{sga3:iv:theorem:4:4:1}{此定理},
可用下面的方法在 \(\widetilde{\mathscr C}\) 中证明同时涉及任意正向极限
和有限逆向极限的陈述 (例如, \textquotedblleft 每个满态射都是普遍有效的\textquotedblright;
参见下文). 先在集合范畴中证明相应陈述, 再逐点推广到预层范畴.
然后利用\SGARefLink{sga3:iv:theorem:4:4:1}{前述定理},
从预层范畴过渡到层范畴. 下文会出现许多使用此方法的例子
(\SGARefLink{sga3:iv:proposition:4:4:3}{4.4.3},
\SGARefLink{sga3:iv:lemma:4:4:6}{4.4.6},
\SGARefLink{sga3:iv:proposition:4:4:9}{4.4.9} 等).}

最后注意, 关于预层范畴的陈述, 形式上是相应的层范畴陈述的推论.
只需取最粗的 (\textquotedblleft 混沌\textquotedblright) 拓扑,
即由
\[
  J(S)=\{S\}
  \qquad\text{对每个 }S\in\operatorname{Ob}\mathscr C;
\]
定义的拓扑; 对此拓扑, 每个函子都是层.

\medskip
\noindent\SGARefTarget{sga3:iv:proposition:4:4:3}\textbf{命题 4.4.3.}
\textit{设 \(\mathscr F=\{F_i\to F\}\) 是层的态射族.
下列条件等价:}
\begin{enumerate}
  \item[(i)]\SGARefTarget{sga3:iv:proposition:4:4:3:item:i:01}
    \textit{\(\mathscr F\) 是满态射族;}
  \item[(ii)]\SGARefTarget{sga3:iv:proposition:4:4:3:item:ii:01}
    \textit{\(\mathscr F\) 是普遍有效满态射族 (\SGARefLink{sga3:iv:remark:1:13}{1.13});}
  \item[(iii)]\SGARefTarget{sga3:iv:proposition:4:4:3:item:iii:01}
    \textit{\(\mathscr F\) 是覆盖族 (\SGARefLink{sga3:iv:definition:4:2:3}{4.2.3});}
  \item[(iv)]\SGARefTarget{sga3:iv:proposition:4:4:3:item:iv:01}
    \textit{\(\mathscr F\) 的层像, 即 \(\mathscr F\) 的预层像
    (\SGARefLink{sga3:iv:definition:4:1:2}{4.1.2}) 的相伴层, 是 \(F\).}
\end{enumerate}

\SGARefLink{sga3:iv:proposition:4:4:3:item:iii:01}{(iii)} 与
\SGARefLink{sga3:iv:proposition:4:4:3:item:iv:01}{(iv)} 的等价性由
\SGARefLink{sga3:iv:corollary:4:3:12}{4.3.12} 得到.
其余等价性将由下面几个引理得出.

\medskip
\noindent\SGARefTarget{sga3:iv:lemma:4:4:4}\textbf{引理 4.4.4.}
\textit{设 \(f:X\to Y\) 是层的单态射, 同时也是满态射. 则 \(f\) 是同构.}

在集合范畴中, \SGARefLink{sga3:iv:lemma:4:4:4}{此引理}首先是显然的.
接着在预层范畴中证明它. 考虑预层
\[
  V:S\longmapsto Y(S)\amalg_{X(S)}Y(S);
\]
这是预层范畴中 \(Y\) 和 \(Y\) 沿 \(X\) 的融合和 (推出).
\footnote{编者注: 证明的其余部分略有修改.}
设 \(i_1,i_2:Y\rightrightarrows V\) 是两个坐标态射.
若 \(X\to Y\) 是预层范畴中的满态射, 则 \(i_1=i_2\).
此时对每个 \(S\), 映射 \(X(S)\to Y(S)\) 都是满射;
由于它也是单射, 故为双射, 因而 \(X\to Y\) 是同构.

最后, 在层范畴 \(\widetilde{\mathscr C}\) 中讨论.
由 \SGARefLink{sga3:iv:theorem:4:4:1:item:ii:01}{4.4.1(ii)},
层 \(Y\) 和 \(Y\) 沿 \(X\) 在 \(\widetilde{\mathscr C}\)
中的融合和是预层 \(V\) 的相伴层 \(Z\). 考虑图
\[
  X\xrightarrow{f}Y
  \overset{i_1}{\underset{i_2}{\rightrightarrows}}
  V\xrightarrow{\tau}Z=a(V).
\]
有 \(i_1f=i_2f\), 因而 \(\tau i_1f=\tau i_2f\),
所以 \(\tau i_1=\tau i_2\), 因为 \(f\) 是
\(\widetilde{\mathscr C}\) 中的满态射.
由\SGARefLink{sga3:iv:lemma:4:4:5}{下一个引理}的
\SGARefLink{sga3:iv:lemma:4:4:5:item:iii:01}{(iii)} 项,
预层 \(V\) 分离, 因此 \(\tau\) 是单态射
(\SGARefLink{sga3:iv:proposition:4:3:11:item:iv:01}{4.3.11(iv)}).
于是 \(i_1=i_2\), 而上面已见这蕴含 \(f:X\to Y\) 是同构.
这就证明了引理~\SGARefLink{sga3:iv:lemma:4:4:4}{4.4.4},
只要再验证 \(V\) 分离即可.

% Translation coverage: entire assigned source component.
% Source lines: 136; translated body lines before this audit comment: 142.
% Math preserved: 43 inline formulas, 9 display blocks.
% Reference targets: 12/12; links: 25/25; labels: 0/0; footnotes: 2/2.
% Unresolved TODOs: 0; compile and visual review delegated to the volume integrator.
